% ============================================================
%  CAPÍTULO 5 — CONCLUSIONES
% ============================================================
\chapter{Conclusiones}
\label{cap:conclusiones}

En este capítulo se presentan las conclusiones del trabajo, organizadas
por objetivo específico, seguidas de los alcances y limitaciones del sistema,
los productos de la investigación y las líneas de trabajo futuro.

% -------------------------------------------------------------
\section{Conclusiones por objetivo}

\textbf{Objetivo 1: Construcción de la base de datos.}
Se construyó una base de datos de imágenes multiespectrales de \ac{MAP}
emuladas enterradas en suelo arenoso con vegetación limitada a \SI{1}{m}
de altura, siguiendo un protocolo reproducible de captura estacionaria mediante
una estructura de soporte de \ac{PVC} diseñada para este fin. La base de datos oficial
comprende seis sesiones de campo (S2--S7) con ocho zonas de mina por sesión
(profundidades de 1, 3, 5 y 7~cm, con dos réplicas por profundidad) y entre
una y cuatro zonas de control (una en S2, cuatro en S3--S7), totalizando
14\,901 ventanas. Las condiciones de iluminación resultaron ser el factor más determinante
para la consistencia de las firmas espectrales entre sesiones: la hora de
captura, la uniformidad de la luz y la humedad del suelo explican en gran
medida las diferencias entre S1 y las sesiones posteriores. Esta observación
motivó el ajuste del protocolo a partir de S2 y la exclusión de la Sesión~1
de la base de datos oficial.

\textbf{Objetivo 2: Procesamiento de imágenes multiespectrales.}
Se implementó un \textit{pipeline} automatizado de calibración radiométrica
mediante \ac{CRP}, alineación de bandas por \ac{SIFT}+RANSAC y extracción de
378 características estadísticas sobre ventanas de $128\times128$ píxeles. Durante la implementación se encontró una limitación no documentada
de la MicaSense RedEdge-MX: la umbralización de Otsu falla en la banda
Red Edge porque la vegetación y el panel de calibración tienen reflectancias
similares a \SI{717}{nm}. Detectar la máscara en la banda Roja y aplicarla
a las cinco bandas resolvió el problema y mejoró la consistencia de la
calibración radiométrica.

\textbf{Objetivo 3: Reconocimiento de patrones.}
El etiquetado preciso del groundtruth ---una ventana es positiva solo si su centro
cae dentro del bounding box de la \ac{MAP}--- produjo 487 ventanas positivas frente
a 14\,414 negativas (desbalance 30:1). El análisis de importancia mostró que las
bandas \ac{NIR} y Red Edge contribuyen de forma clara a la discriminación: de las
70 características del modelo final, 36 (la mitad) involucran \ac{NIR} o Red Edge. Tras comparar \ac{RF},
\ac{SVM}, XGBoost y EasyEnsemble, se eligió \textbf{EasyEnsemble~x10 XGBoost}
(submuestreo 1:1 para corregir el desbalance 30:1) con las \textbf{top-70
características} de mayor ganancia de XGBoost, alcanzando
\textbf{AUC~=~0.947~$\pm$~0.007} y sensibilidad del 82.3\,\% a umbral~0.5. Frente a
la línea base (\ac{RF}, 378 características, AUC~=~0.915), la ganancia total es de
$+$0.032 AUC.

\textbf{Objetivo 4: Evaluación del sistema.}
El modelo detectó minas en las cuatro profundidades evaluadas. La sensibilidad
disminuye con la profundidad ---del 87.5\,\% a \SI{1}{cm} al 70.7\,\% a
\SI{7}{cm}---, coherente con que la señal proviene de la perturbación superficial
del suelo y del estrés vegetal: a mayor profundidad, menor alteración detectable.
Aun así, el \ac{AUC} se mantiene alto en todas las profundidades (0.91--0.94), muy
por encima del azar. Respecto a los falsos positivos, el \textbf{93\,\%} no se
asocia a los objetos de interferencia (botella, lata, piedra), sino al borde
difuso del suelo perturbado alrededor de la mina; el clasificador no confunde
sistemáticamente esos objetos con la \ac{MAP}.

% -------------------------------------------------------------
\section{Alcances y limitaciones}

\subsection{Alcances}

\begin{itemize}
  \item El sistema alcanzó AUC~=~0.947 con Sens(0.5)~=~82.3\,\% (umbral 0.5)
    sobre \ac{MAP} emuladas en suelo arenoso con vegetación hasta \SI{1}{m},
    usando EasyEnsemble~x10 XGBoost con las top-70 características seleccionadas
    por ganancia de XGBoost bajo validación LOSO.
  \item El \textit{pipeline} cubre todo el proceso sin intervención manual,
    desde la captura hasta el mapa de probabilidades. Esto lo hace
    directamente aplicable en futuras campañas de campo.
  \item Las 378 características pueden reducirse a las top-70 sin pérdida de
    rendimiento (incluso ligeramente superior a usar las 378). El vector de
    entrada se reduce a menos de una quinta parte.
  \item El análisis de importancia mostró que las bandas Red Edge y \ac{NIR}
    contribuyen a la detección: 36 de las 70 características del modelo final las
    involucran.
  \item La validación \ac{LOSO} evita que ventanas de la misma sesión aparezcan en
    entrenamiento y prueba al mismo tiempo, lo que haría los resultados
    irrealmente optimistas. Además, una validación que deja fuera cada
    emplazamiento físico (\textit{leave-one-zone-out}, AUC~=~0.933 frente a 0.947)
    indica que el modelo generaliza a sitios no vistos bajo condiciones ya
    observadas.
\end{itemize}

\subsection{Limitaciones}

\begin{itemize}
  \item La base de datos usa minas emuladas en suelo arenoso bajo condiciones
    controladas. No se ha probado el sistema con minas reales ni en terrenos
    con mayor heterogeneidad de suelo o vegetación.
  \item El sistema solo distingue entre mina y no-mina. Aunque el análisis
    de FP incluyó objetos de interferencia (botella, lata, piedra),
    clasificarlos por separado no fue parte del trabajo.
  \item El modelo responde a la firma espectral de estrés vegetal y
    perturbación del suelo asociadas a la mina. Zonas con vegetación
    naturalmente deteriorada o suelo removido por causas ajenas pueden
    presentar firmas similares, contribuyendo a las falsas alarmas.
  \item Las condiciones extremas de iluminación pueden afectar las firmas
    espectrales, como evidenció la Sesión~1 (HR~=~90.5\,\%, ángulo solar
    rasante). Dentro del rango del protocolo (S2--S7), el análisis por
    sesión no mostró una correlación fuerte entre la hora de captura y el
    AUC, aunque el sistema requiere un protocolo de captura estricto para
    evitar condiciones atípicas.
  \item El paralaje entre los cinco lentes de la MicaSense a \SI{1}{m}
    genera un desalineamiento residual que la homografía global no puede
    corregir completamente sobre vegetación tridimensional.
  \item El ratio positivos:negativos de 1:30 limita la cantidad de ventanas de
    mina disponibles; más sesiones o más minas por sesión mejorarían la robustez
    del modelo.
\end{itemize}

% -------------------------------------------------------------
\section{Productos del trabajo de investigación}

A partir del desarrollo de este trabajo de grado se obtuvieron los
siguientes productos:

\begin{itemize}
  \item \textbf{Documento de trabajo de grado.} El presente documento,
    que constituye la memoria técnica completa del sistema desarrollado.

  \item \textbf{Base de datos de imágenes multiespectrales.} Base de datos de
    imágenes de \ac{MAP} emuladas enterradas en suelo arenoso con vegetación
    limitada, adquiridas con la cámara MicaSense RedEdge-MX
    a \SI{1}{m} de altura en el campus de la Universidad del Valle.
    La base de datos incluye ocho zonas de mina por sesión (profundidades de
    1, 3, 5 y 7~cm) y cuatro zonas de control, con seis sesiones de
    campo bajo distintas condiciones ambientales.

  \item \textbf{Código fuente del \textit{pipeline} de procesamiento.}
    Implementación completa en Python del sistema de preprocesamiento
    radiométrico y geométrico, extracción de características espectrales
    y clasificación con validación \ac{LOSO}. 
\end{itemize}

% -------------------------------------------------------------
\section{Trabajos futuros}
\label{sec:trabajos_futuros}

Con base en los resultados y limitaciones identificadas, se proponen las
siguientes líneas de trabajo futuro:

\begin{enumerate}
  \item \textbf{Clasificación multiclase y reducción de falsos positivos.}
    Aprovechar las anotaciones de cuatro clases (MAP, Botella, Lata, Piedra)
    ya disponibles en la base de datos para desarrollar un clasificador capaz de
    distinguir entre \ac{MAP}, objetos inofensivos y suelo limpio, reduciendo
    así la tasa de falsos positivos del sistema actual. El análisis espectral
    por clase (Figura~\ref{fig:firmas_espectrales}) muestra que las cuatro
    clases presentan patrones parcialmente diferenciables, particularmente en
    las bandas NIR y Red Edge, lo que respalda la factibilidad de este enfoque.

  \item \textbf{Montaje en UAV.} Escalar el sistema a un \ac{UAV} para
    validar la operación en condiciones de vuelo real, considerando el
    movimiento de la plataforma y su efecto sobre la alineación de bandas.
    Los resultados del trabajo previo del grupo \ac{PSI} \cite{tenorio2024}
    en la misma plataforma constituyen una referencia directa para este paso.

  \item \textbf{Fusión de sensores.} Integrar la información multiespectral
    con datos termográficos del enfoque previo del grupo \ac{PSI}
    \cite{tenorio2024} para explorar si la fusión de ambas modalidades
    ---señal crónica multiespectral y señal aguda térmica--- mejora la
    tasa de detección y reduce los falsos positivos.

  \item \textbf{Análisis de estabilidad espectral temporal.} Estudiar cómo
    evolucionan las firmas espectrales sobre las \ac{MAP} a lo largo de
    semanas y meses, para identificar la ventana temporal óptima de detección
    basada en estrés vegetal crónico y determinar si el efecto aumenta,
    se estabiliza o decrece con el tiempo.

  \item \textbf{Generalización a terrenos heterogéneos.} Validar el sistema
    en terrenos con mayor densidad vegetal, vegetación más alta, o suelos
    con diferente textura y composición, para establecer los límites
    operativos del enfoque multiespectral.

  \item \textbf{Incertidumbre del sensor y sesgo del NDRE.} Incorporar la
    incertidumbre de la calibración radiométrica en el análisis de resultados,
    en particular el sesgo residual del \ac{NDRE} documentado en este trabajo,
    para cuantificar su efecto sobre las características más discriminativas.
\end{enumerate}
